\chapter{Limitations}
\label{ch:limitations}

This chapter states the limitations of the evidence and implementation
underlying Chapters~\ref{ch:activities}--\ref{ch:interpretation}, organised
by category rather than by project phase. It does not re-run the results
chapter: each item below is stated once, concisely, with a pointer to
where it was established, and the consequence for what the report can and
cannot claim.

\section{Dataset and Experimental-Design Limitations}
\label{sec:lim-dataset}

Every result in this report ultimately rests on 48 physical specimens
tested across only two experimental campaigns, within which steel-mesh
configuration, chloride concentration, and terminal ageing duration
co-vary almost perfectly (Section~\ref{sec:confounding}). Structural
supervision is available at a single terminal timepoint per specimen, with
no repeated or intermediate structural measurement and no observed
failure-event time (Section~\ref{sec:structural-supervision}). The
photographic record covers external surface appearance only; nothing in
the dataset directly measures internal or subsurface condition at any
non-terminal timepoint. Treatment groups are uneven in size, from two
specimens in the smallest protocol to eighteen in the largest
(Table~\ref{tab:campaign-design}). Together, these properties bound every
downstream claim about structural condition, degradation, or transfer
across campaigns or treatments to what a two-campaign, 48-specimen design
can support — a limitation of the physical experiment this project
inherited, not of the modelling choices made within it.

\section{Label and Target Limitations}
\label{sec:lim-labels}

Visible-corrosion labels are manually assigned quantities
(Section~\ref{sec:data-modalities}), not an independent instrumental
ground truth, and two independent findings show that a benchmark built
against them can be label-adjacent rather than purely image-driven: the
classical baseline's peak-rust prediction is dominated by a rust-mask
feature closely aligned with the label (Section~\ref{sec:phase1}), and the
robustness-quantification phase's total-rust benchmark uses a feature
correlated with its label at Spearman $\rho=0.9996$
(Section~\ref{sec:phase4a}). The category schema used for the two
severity targets differs between dataset builds — five ordinal levels
historically, four currently (Table~\ref{tab:category-schema}) — a
genuine version difference in the underlying workbook, not a data-quality
error, but one that prevents any five-class and four-class classification
result from being compared directly. Proxy-RUL thresholds (wire-loss
percentages, a load ratio, a composite health-index cutoff, or the refined
pipeline's single-series thresholds) are engineering screening choices,
not values derived from an independently validated failure criterion
(Section~\ref{sec:proxy-rul}), and no ground-truth remaining-useful-life
value exists anywhere in the dataset against which any of them could be
checked.

\section{Generalisation Limitations}
\label{sec:lim-generalisation}

A grouped specimen holdout answers a narrower question than transfer to a
new treatment or campaign (Section~\ref{sec:grouped-holdout}), and
leave-one-campaign-out in this dataset is better understood as a severe
extrapolation stress test than as an estimate of ordinary future
performance, given how completely campaign identity is entangled with the
principal design variables (Section~\ref{sec:loco}). Treatment-level
(leave-one-treatment-out) behaviour for visible-corrosion targets is not a
settled project-wide property: it is a genuine failure mode in the
interpretable-feature phase and is not reproduced in the
robustness-quantification phase's own treatment-holdout evaluation
(Section~\ref{sec:results-visible}), and the report does not know why.
Within-mesh subgroup analyses are based on only twenty-four specimens per
group (Section~\ref{sec:phase4b}), which limits how much confidence any
single subgroup result — positive or negative — can carry on its own.

\section{Structural-Modelling Limitations}
\label{sec:lim-structural}

Every structural claim in this report is built on 48 labelled specimens
(Section~\ref{sec:lim-dataset}); no amount of downstream modelling
sophistication changes this. The interpretable-feature phase's structural
models are mixed-input — metadata, manually assigned corrosion labels, and
image descriptors together, not image-only
(Section~\ref{sec:phase3-structural}) — while the robustness-quantification
phase's final, robustness-weighted selections use metadata alone
(Section~\ref{sec:results-structural}); no phase in this project has
produced a robustness-weighted structural model shown to rely on image
input. The degradation stage's structural estimates remain in-sample
refits rather than out-of-fold predictions (Section~\ref{sec:eng-contracts}),
so no validated longitudinal structural state exists at any non-terminal
week for any specimen — every such value in the dataset is model-derived,
not measured (Table~\ref{tab:quantity-availability}). As established in
Chapter~\ref{ch:interpretation}, this bounds what can be inferred to
predictive identifiability from this specific dataset; it does not bound,
and this report does not claim to bound, whether a physical relationship
between surface corrosion and structural capacity exists.

\section{Prognostic / Proxy-RUL Limitations}
\label{sec:lim-rul}

No true remaining-useful-life supervision exists in this dataset
(Section~\ref{sec:proxy-rul}), and this report's use of "proxy-RUL"
throughout is a deliberate, consistently applied qualification of that
fact. Every threshold used to derive a proxy-RUL or threshold-status value
is a screening assumption, not a validated failure criterion
(Section~\ref{sec:lim-labels}). Many specimens are already past the
relevant threshold at their first observation in the refined pipeline's
own accounting, and that same pipeline's own output shows \emph{zero
future} threshold crossings across every tested case — not zero crossings
overall, since baseline- and during-observation-crossed cases are common
(Section~\ref{sec:results-degradation}). No out-of-fold structural
trajectory feeds either generation of the degradation model
(Section~\ref{sec:eng-contracts}), and no validation exists against an
observed future failure or a repeated structural measurement of any
specimen. None of this makes the pipeline useless: as
Chapter~\ref{ch:interpretation} establishes, its value is methodological
and diagnostic — it clarifies exactly what data a validated version would
require — and that value is independent of, and unaffected by, its not
being a validated prognostic method.

\section{Current Classification-Branch Limitations}
\label{sec:lim-classification}

The four-class severity schema is severely imbalanced, with the majority
class accounting for over four-fifths of images
(Table~\ref{tab:data-modalities}). No classifier has been trained,
compared, or evaluated on this data at the time of writing
(Section~\ref{sec:phase5}), so no accuracy, macro-F1, or confusion-matrix
result exists to limit. Uniform six-fold augmentation increases the
absolute amount of training data without changing the relative class
proportions, and does not by itself address the imbalance
(Section~\ref{sec:phase5}). Because no model has been trained, performance
on held-out or external data is entirely unknown; this report does not
speculate about the accuracy a ResNet-50 or Vision Transformer
architecture might achieve, since neither has been implemented for this
task.

\section{Software and Reproducibility Limitations}
\label{sec:lim-software}

Stated briefly here and detailed in Chapter~\ref{ch:engineering}, to which
this section defers rather than repeats: a live YAML boolean-parsing bug
and a broader "1"/"first"~$\rightarrow$~"main\_first" string-substitution
pattern currently prevent an unmodified rerun of several \texttt{main\_4}
pipeline stages (Section~\ref{sec:eng-schema}); \texttt{main\_4} has no
exported dependency specification (Section~\ref{sec:eng-evolution}); and
two of the five artifact-contract improvements the project's own
documentation recommends remain unimplemented
(Section~\ref{sec:eng-contracts}). These are reproducibility limitations
of the current source tree, not limitations of the scientific evidence
itself: the saved scientific-output artifacts from the affected
\texttt{main}--\texttt{main\_4} pipeline generations that underpin the
regression, structural, degradation, and robustness results reported here
were generated before the earliest observed source-corruption timestamp
(Section~\ref{sec:eng-schema}) and remain valid on that basis. The later
\texttt{augmentation} branch, which underpins the four-class
classification data preparation, is independently confirmed unaffected by
this corruption (Section~\ref{sec:eng-schema}) and is not part of this
claim.
